\documentclass[a4paper,12pt]{article}
\usepackage[utf8]{inputenc}
\usepackage[T1]{fontenc}
\usepackage[english]{babel}
\usepackage{geometry}
\geometry{left=2.5cm,right=2.5cm,top=2.5cm,bottom=2.5cm}

\usepackage{lmodern}
\usepackage{xcolor}
\definecolor{titleblue}{RGB}{0,51,140}

\usepackage{hyperref}

\pagestyle{empty}

\begin{document}

\begin{flushleft}
    \textbf{Ismail AISSAOUI} \\
    State Engineer in Artificial Intelligence \\
    ismail.aissaoui.pro@gmail.com \\
    +213 660 70 77 96 \\
    \href{https://ismail-aissaoui.vercel.app}{ismail-aissaoui.vercel.app}
\end{flushleft}

\begin{flushright}
    To the PhD Selection Committee \\
    \textbf{EDT Program - CEA List / IRIT} \\
    Attn: Pr. Brahim Hamid and Dr. Marcos Didonet Del Fabro \\
    \medskip
    May 1, 2026
\end{flushright}

\bigskip

\noindent \textbf{Subject: Application for PhD: AI-Assisted Model-Driven Security Engineering for Digital Twins (Rank 5 - PC3)}

\bigskip

Dear Pr. Hamid and Dr. Del Fabro,

As a State Engineer in Artificial Intelligence from ESI-SBA, I am writing to express my enthusiastic application for the PhD position in "AI-Assisted Model-Driven Security Engineering for Digital Twins" within the Engineering Digital Twin (EDT) national program. This application is motivated by the alignment between my background in GNN-based security and the AI-assisted modeling objectives of your research group at CEA/IRIT.

Currently, I am completing my thesis at Sonatrach, where I have developed a **Physics-Informed Digital Twin** for industrial monitoring. Beyond the core physics-aware modeling, I have a strong background in **Relational Intelligence** and **Graph Neural Networks (GNNs)**, which I previously applied to cybersecurity tasks such as structural anomaly detection in complex networks. I am deeply interested in how AI can move beyond simple code assistance to become a systematic partner in the Model-Driven Engineering (MDE) of secure digital twins.

I am particularly excited by your proposal to close the loop between design-time security models and runtime execution feedback using **round-trip engineering**. My experience in deploying real-time AI models on industrial edge devices has taught me the critical importance of synchronizing architectural models with live vulnerability scans and attack simulations. I am eager to apply my skills in generative AI and relational modeling to develop an AI-assisted environment that helps designers inject security mitigations directly into the Asset Administration Shell (AAS) specifications.

My background in GNN-based security and my experience with production-grade PyTorch and Model-Driven Engineering (MDE) provide me with the necessary tools to contribute to the **Artemis platform** and the broader EDT ecosystem. I am a highly autonomous researcher, comfortable bridging the gap between automated security analysis and complex architectural modeling.

Thank you for your time and consideration. I look forward to the possibility of discussing how my experience in AI-assisted modeling and cybersecurity can contribute to the success of your research group at CEA List and IRIT.

Sincerely,

\bigskip

\textbf{Ismail AISSAOUI}

\end{document}
